\subsection{The topology of locally parameterized outer faces}

Pseudoisotopy obstructions to families of cone structures also appear
in Ayala--Francis--Tanaka's comparison between naive smooth families
and conically smooth stratified spaces.  In their category, the
automorphism space of a compact link is equivalent to that of its
open cone \cite[arXiv version~6, pp.~3--4; Theorem~4.3.1,
p.~46]{AyalaFrancisTanaka2017v6}.
The family category and the topology on the space of faces are
therefore part of the question.
Nariman's model for a three-manifold homeomorphism-group classifying space uses a
quotient of locally flat embeddings into an infinite-dimensional
half-space \cite[Section~7, pp.~13--14, especially Definition~7.1]
{Nariman2025}.  Here a fixed finite-dimensional cone and its shell
determine the stabilizer.
Salvatore--Turchin's smoothing-theory deloopings concern spaces of
disc embeddings and diffeomorphisms.  For $1\le m\le n$, their
Theorem~B excludes ambient dimension four and restricts to
invertible components in codimension two
\cite[pp.~2--3]{SalvatoreTurchin2026}.
The argument below uses compact-open locally flat embeddings,
supported isotopy extension, and a contraction of the ambient
cone homeomorphism group.

Let $B$ be a closed connected PL manifold of dimension $d\geq4$.
Write $X=C_2B$ for the cone with radial interval $[0,2]$, let $o$ be
its vertex, and put $C=C_1B$ and $Y_*=B\times\{1\}$.
A strict cone neighborhood $(D,Y)$ has a closed cone image $D\subset X$,
an outer face $Y\subset B\times(0,2)$, a homeomorphism $CY\to D$
fixing the base and taking the vertex to $o$, an open deleted-face
part $D\setminus Y$, and an exterior collar of $Y$.
The cone chart and exterior collar give an interior and an exterior
collar, so $Y$ is two-sided and locally flat.
The cone chart is required to exist but is not retained in the object.

The face determines its cone image: $D\setminus Y$ is the connected
component of $o$ in $X\setminus Y$, since it is connected and both
open and closed there, and $D=\overline{D\setminus Y}$.
We give outer faces the subspace topology in
\[
 \operatorname{Emb}_{\mathrm{lf}}(B,B\times(0,2))
       /\operatorname{Homeo}(B),
\]
where embeddings have the compact-open topology and the action is by
precomposition.  This topology is specified before any choice of
parameterization is forgotten.

A local joint face witness over a parameter neighborhood $U$ is a
homeomorphism
\[
 B\times U\longrightarrow
       \{(p,y):p\in U,\ y\in Y_p\}
\]
over $U$, with jointly continuous forward and inverse maps.
In a semialgebraic family these maps are semialgebraic.
A finite semialgebraic simplex filling consists of a finite
triangulated parameter disk, semialgebraic cone and face families
on its closed simplices, and a joint semialgebraic face witness
on each simplex.  Cone and face images agree on common parameter
faces and equal the prescribed loop on the boundary of the disk.
The witness parameterizations need not agree and are not retained.
Over a real closed field the cone charts, collar witnesses, and
their inverses are also required to be semialgebraic; ``compact''
means closed and bounded.  This finite convention imposes no
retained radial ruling and invokes no topological finite-subcover
property.

Put
\[
 P_{\mathrm{TOP}}(B)=
 \operatorname{Homeo}(B\times[1,2]\ \mathrm{rel}\ B\times\{2\}).
\]
Interval reversal identifies this with the usual pseudoisotopy group.

\begin{theorem}\label{thm:locally-parameterized-outer-faces}
The component $\mathcal W$ of the standard cone in this outer-face
space has weak homotopy type $BP_{\mathrm{TOP}}(B)$.
For $B=T^6$ there is a semialgebraic loop of strict cones that does
not bound a real definable family with local joint face witnesses.
The loop can be defined over $\mathbb Q$, and over every real closed
field it has no finite semialgebraic simplex filling with the
specified witnesses.

If a parameterization $B\to Y$ is retained, the component of the
standard parameterized face instead has weak homotopy type
\[
 B\operatorname{Homeo}
       (B\times[1,2]\ \mathrm{rel}\ B\times\{1,2\}).
\]
\end{theorem}

\begin{proof}
\emph{The quotient metric and finite-parameter face bundles.}
Let $E$ be the space of embeddings occurring as strict outer faces,
and let $U=E/\operatorname{Homeo}(B)$.
Fix a compatible metric $d_X$ on the compact cone.
The compact-open topology on $E$ is the topology of
$d_\infty(e,f)=\sup_{b\in B}d_X(e(b),f(b))$.
Precomposition acts by isometries, and
\[
 \delta([e],[f])=
 \inf_{a\in\operatorname{Homeo}(B)}
       \sup_{b\in B}d_X(e(b),f(a(b)))
\]
induces the quotient topology.
Indeed, inverse images of $\delta$-balls are unions of uniform balls;
conversely, an open saturated set containing $e$ contains a uniform
ball about $e$, whose quotient contains the corresponding
$\delta$-ball.  This is a metric: the Hausdorff distance between
$e(B)$ and $f(B)$ is bounded by $\delta$, and equal images differ by
a homeomorphism of $B$.

Let $P$ be a compact finite-dimensional metric space and let
$p\mapsto Y_p$ be a continuous map $P\to U$.  Its total face space
\[
 Z=\{(p,x)\in P\times X:x\in Y_p\}
\]
is closed in $P\times X$.  If $(p_i,x_i)\to(p,x)$ with
$x_i\in Y_{p_i}$, quotient-metric continuity implies Hausdorff
convergence to $Y_p$, and hence $x\in Y_p$.
Thus $Z$ is compact and its projection to $P$ is proper.
For any $p$ and $\varepsilon>0$, sufficiently close $q$ satisfy
$d_P(p,q)<\varepsilon/2$ and $\delta(Y_p,Y_q)<\varepsilon/2$.
The latter inequality supplies a homeomorphism between the two
fibres moving every point by less than $\varepsilon/2$ in $X$,
and therefore by less than $\varepsilon$ in $Z$ with its product
metric.  The projection is completely regular.

The completely regular mapping theorem
\cite[Proposition~7.1 and preceding definition, p.~605]{ChapmanFerry1979}
now applies.  Its base is locally finite-dimensional, its projection
is proper, and the homeomorphism group of the closed manifold fibre
is locally contractible by
\cite[Corollary~1.1, p.~63]{EdwardsKirby1971}.
Consequently $Z\to P$ is a locally trivial $B$-bundle.
In particular, a quotient family over a disk has a joint face
parameterization.  Over $D^k\times I$ a trivialization can be made
to extend any prescribed bottom parameterization: its discrepancy
at $t=0$ is a continuous map $D^k\to\operatorname{Homeo}(B)$,
and precomposition by that map, held constant in $t$, gives the
required correction.  Thus $E\to U$ has the Serre lifting property.
Only finite-dimensional compact parameter bases have been used.

\emph{The ambient action.}
Consider
\[
 G=\operatorname{Homeo}
       (X\ \mathrm{rel}\ \{o\}\cup B\times\{2\}).
\]
This group is contractible.  If $s_t:X\to C_{2t}B$ denotes radial
scaling, set, for $t>0$,
\[
 A_t(g)(x)=
 \begin{cases}
  s_tg s_t^{-1}(x),&x\in C_{2t}B,\\
  x,&x\notin C_{2t}B,
 \end{cases}
 \qquad A_0(g)=1.
\]
The formulas agree on the joining face because $g$ fixes the outer
base.  The inverse is $A_t(g^{-1})$.
At $t=0$ the displacement of both maps is bounded, uniformly in
$g$, by the diameter of $C_{2t}B$, which tends to zero in a conical
metric.  This proves continuity of the Alexander contraction,
including its dependence on $g$.

We use the compact-family extension theorem of
\cite[Theorem~1, pp.~1247--1248; covering-isotopy theorem,
p.~1250]{Chernavskii1969}.
In that theorem a compact family is jointly continuous and each
fibre embedding is locally flat; a locally flat total parameter
track is not an additional hypothesis.
Here the source is compact, the ambient manifold $B\times(0,2)$
has dimension $d+1\geq5$, and the codimension is one, hence differs
from two.  The local complementary simple-connectivity condition
in the small-extension theorem holds in a local flat chart.
Its displacement control function can vanish off a neighborhood
of the embedded source.  The extension of a compact parameter
family can therefore be supported in a compact annulus containing
its trace, and extends by the identity to an element of $G$.

The small-extension part also shows that every ambient orbit in
$E$ and in $U$ is open, and that its orbit map from $G$ is open.
For a prescribed identity neighborhood in $G$, choose the ambient
push small enough in both the forward and inverse uniform metrics.
Uniformly close parameterized embeddings then lie in its orbit
image; after reparameterization the same holds in the quotient.
Every orbit is connected, since $G$ is contractible, and so the
orbits are exactly the connected components.
Write $E_0$ and $U_0=\mathcal W$ for the standard components.

The map $G\to E_0$, $g\mapsto g|_{Y_*}$, is a Serre fibration.
Indeed, a family over $D^k\times I$ extends by the compact-family
theorem to a continuous family of ambient maps.  If $g(z,t)$ is
one such lift and $g_0(z)$ is the prescribed bottom lift, then
\[
 \widetilde g(z,t)=g(z,t)g(z,0)^{-1}g_0(z)
\]
has the prescribed initial value and still induces the same face
parameterization, since the right-hand correction fixes $Y_*$
pointwise.  The lifting property of $E\to U$ restricts to
$E_0\to U_0$, since lifted paths remain in their component.
It follows that $G\to U_0$ is a Serre fibration as well.

\emph{The two stabilizers.}
Let $H\subset G$ preserve $Y_*$ setwise, and let $H^{\mathrm{pt}}$
fix it pointwise.  The open orbit maps identify
$U_0=G/H$ and $E_0=G/H^{\mathrm{pt}}$ as topological quotients.
Put
\[
 A=\operatorname{Homeo}(C\ \mathrm{rel}\ \{o\}\cup Y_*),
 \qquad
 Q=\operatorname{Homeo}
       (B\times[1,2]\ \mathrm{rel}\ B\times\{1,2\}).
\]
Restriction to the outer shell gives a split group homomorphism
$H\to P_{\mathrm{TOP}}(B)$.  Its section extends the inner-end
boundary map conically over $C$, and its kernel is $A$.
Thus $H\cong A\rtimes P_{\mathrm{TOP}}(B)$, with a product
decomposition of the underlying spaces.  The group $A$ has the
same Alexander contraction.  For the pointwise stabilizer,
independent restrictions to the cone and shell give
$H^{\mathrm{pt}}\cong A\times Q$.

The principal Serre fibrations with contractible total group now
give
\[
 U_0\simeq_w BH\simeq_w BP_{\mathrm{TOP}}(B),
 \qquad
 E_0\simeq_w BH^{\mathrm{pt}}\simeq_w BQ.
\]
One may obtain the classifying-space conclusion directly on
singular simplicial sets: every simplex in the quotient lifts,
and two lifts differ by a unique continuous simplex in the
stabilizer.  The quotient of the contractible singular total
group therefore has the stabilizer's classifying weak type.
This uses no global section over the infinite-dimensional
quotient.

\emph{A rational semialgebraic essential loop.}
For $B=T^6$, the classical torus component calculation gives
\[
 \pi_0P_{\mathrm{PL}}(B)
       \cong\mathbb F_2[\mathbb Z^6]/\mathbb F_2[0]\ne0,
\]
and these classes survive in $P_{\mathrm{TOP}}(B)$, by
Theorem~\ref{result-3} and its unstable category comparison
\cite[Theorems~6.1--6.2, pp.~41--43]{BurgheleaLashof1974II}.
Choose a nonzero PL representative $p$ on the shell, fixed at its
outer end.  Use a rational finite PL model of the torus and a
rational PL product identification of the shell.

The representative may be chosen with rational PL coefficients.
To see this, choose finite source and target triangulations on
which $p$ is simplicial, refining a fixed rational triangulation
of the shell and agreeing on the fixed end.  Perturb their
vertices to rational points in their carriers, identically on
the fixed end.  Sufficiently small carrier-preserving perturbations
of this finite triangulation preserve its simplex incidences,
nondegeneracy, and coverage; the straight vertex motions therefore
give continuous families of PL homeomorphisms
$\alpha_s,\beta_s$ of the shell.
They agree on the fixed end, so $\beta_s p\alpha_s^{-1}$ is a
path in the topological pseudoisotopy group and its final member
has rational affine pieces.  Its TOP component is therefore the
same nonzero component as that of $p$.  Replace $p$ by that member.

Cone the free-end map of $p$ to obtain a rational PL element
$h\in H$, and set
\[
 g_t=A_t(h),\qquad C_t=g_t(C),\qquad Y_t=g_t(Y_*).
\]
The images form a loop because $g_0=1$ and $g_1=h\in H$.
The family is semialgebraic over $\mathbb Q$: on an affine piece
of $h$, the scaled expression $t\,h(x/t)$ is affine in $(x,t)$,
and its piece conditions are homogeneous linear inequalities.
The inverse has the same finite form using $h^{-1}$.
Transporting the standard cone chart and collar supplies joint
semialgebraic witnesses on the parameter interval.
If $k=h|_{Y_*}$, precompose the terminal face parameterizations
by $k^{-1}$ near $t=1$.  Their limiting value then equals the
initial parameterization, giving a joint local witness across
the identified endpoints of the circle.

The boundary homomorphism
$\pi_1(U_0)\to\pi_0(H)$ sends this loop to $[h]$, up to inverse.
Shell restriction sends it to the chosen nonzero class $[p]$.
Hence the loop is nonzero in $\pi_1(U_0)$.

\emph{Real fillings.}
A real family with local joint face witnesses gives a continuous
map to $U$: compactness of $B$ turns each total-space witness into
a continuous map to $E$, and equality of images glues the quotient
maps on overlaps.  A disk filling would therefore nullhomotope the
essential loop.  This excludes the asserted real definable fillings;
the argument applies to every continuous filling with these
witnesses.  No conclusion about Hausdorff-only image families is
being used.

\emph{Finite simplex fillings over real closed fields.}
Suppose such a filling exists over a real closed field $F$.
Its finitely many simplex families already supply finite graph
formulas for the joint face witnesses and their inverses.
To retain strictness under transfer, we first produce finitely
many graph formulas for cone and collar witnesses.

On each closed parameter simplex $\sigma$, apply the finite-subset
form of semialgebraic triviality to $X\times\sigma\to\sigma$,
compatible with the total cone and face families,
$\{o\}\times\sigma$, and $(B\times\{2\})\times\sigma$
\cite{Hardt1980}\cite[Theorem~5.46]{BasuPollackRoy2006}.
There is a finite semialgebraic partition into pieces $T_j$ with
ambient-pair trivializations.  Choose one parameter in each
nonempty piece and one semialgebraic cone chart and exterior
collar for its fibre.  Transport these along the trivialization.
If $a_z$ is its ambient map and $\phi_j:CY_j\to C_j$ the chosen
cone chart, the transported chart is
\[
 [y,r]\longmapsto a_z\phi_j([a_z^{-1}(y),r]).
\]
It fixes its base and vertex; the same conjugation transports
the collar.  These joint maps and their inverses are semialgebraic
on their indicated pieces.  They need not agree across pieces,
since they witness strictness and are not part of the object.
The construction uses a finite partition theorem, not a finite
subcover of a non-Archimedean topological space.

Fix finite polynomial formulas for all these sets and maps,
and collect their coefficients into one finite tuple $c$.
The parameter complex may be replaced by the rational realization
of its finite abstract triangulation, retaining the graph of its
boundary identification as part of this finite datum.
For these fixed formulas, being graphs of mutually inverse maps,
continuity, preservation of base coordinates, closedness and
boundedness of cone fibres, openness of their deleted-face parts,
the cone and collar identities, and equality of images on common
faces and on the boundary are first-order conditions.
Continuity is the epsilon--delta formula with the domain
predicates included.  The cone domain can be represented by
\[
 \{(ry,r):y\in Y,\ 0<r\leq1\}\cup\{0\},
\]
so no quotient sort is required.  In particular, no unrestricted
quantification over all semialgebraic chart formulas occurs.

All coefficients of the fixed torus, cone, and loop are rational.
Thus the existence of this finite-format filling is a sentence
$\exists c\,\Phi(c)$ over $\mathbb Q$.
Transfer for real closed fields makes it true over $\mathbb R$
\cite[Theorems~2.80--2.81]{BasuPollackRoy2006}.
The resulting real datum gives a continuous quotient map on each
closed parameter simplex.  Their image agreement and the finite
closed-simplex pasting lemma give a continuous disk in $U$ with
the prescribed boundary loop, contrary to the preceding paragraph.
This proves the field-independent non-filling assertion.
\end{proof}
